\documentclass[11pt,a4paper]{article}
\usepackage[utf8]{inputenc}
\usepackage[T1]{fontenc}
\usepackage{amsmath,amsfonts,amssymb}
\usepackage{graphicx}
\usepackage[hypertexnames=false]{hyperref}
\usepackage{listings}
\usepackage{color}
\usepackage{booktabs}
\usepackage{float}
\usepackage{geometry}
\geometry{margin=1in}
\definecolor{zenblue}{RGB}{41,121,255}
\definecolor{codegray}{RGB}{245,245,245}
\hypersetup{colorlinks=true,linkcolor=zenblue,urlcolor=zenblue,citecolor=zenblue}
\lstset{basicstyle=\ttfamily\small,backgroundcolor=\color{codegray},frame=none,
        breaklines=true,columns=fullflexible,xleftmargin=1em}

\title{\textbf{MoDE v3: Mixture of Diverse Experts}\\
\large Complexity-Aware Routing over Frozen, Device-Pinned Experts\\
Harvested from Open Models}
\author{Zen LM Research, Hanzo AI\\
\texttt{research@hanzo.ai}}
\date{2026}

\begin{document}
\maketitle

\begin{abstract}
We present MoDE v3 (Mixture of Diverse Experts), the architecture specified for the zen
model family. It is not deployed: the zen SKUs serve today from upstream open models while
MoDE is built, and zen5-max is intended to be the first model that runs it. MoDE
inverts the economics of frontier training: rather than train experts, it \emph{harvests}
them. Expert modules are extracted from the largest open-weight models available---six
architecturally distinct families---and served frozen. The only trained parameters are a
complexity estimator, a cross-architecture alignment layer, and a per-tier router:
$\sim$394M, or 0.013\% of the assembled $\sim$3.1T pool. A hierarchical router estimates
task difficulty from a prompt's first tokens and activates only the tier that task
warrants, so a greeting costs 0.8B active parameters and a proof costs 100B+.

The second half of this paper addresses what a 3.1T pool actually requires to run:
\emph{placement}. Because harvested experts are frozen, a pinned expert is a pure
function that is permanently resident on one device---it never syncs, holds no optimizer
state, and never reloads. We show that this makes expert placement a static value rather
than a scheduling problem, and that it extends the kernel-fusion algebra by exactly one
class. Where the existing algebra classes an operation as \textsc{Map} (composes) or
\textsc{Reduce} (fences within a device), MoDE adds \textsc{Route} (fences \emph{across}
devices). The same rule that makes fusion decidable then makes placement decidable.
There is no evaluation: no MoDE model has been built, and the premise the architecture rests
on---that a trained projection makes another family's frozen FFN useful---is itself untested.
We publish no numbers we have not measured.
\end{abstract}

\section*{Lineage}

MoDE has three versions, and they differ in where the experts come from. v3 completes the
line.

\begin{table}[H]
\centering
\caption{The MoDE line. v1 and v2 established the concepts; v3 completes them.}
\begin{tabular}{llll}
\toprule
& \textbf{Name} & \textbf{Where experts come from} & \textbf{Status} \\
\midrule
v1 & Mixture of \emph{Distilled} Experts & Upcycled from our own dense checkpoints & Superseded \\
v2 & Mixture of \emph{Diverse} Experts    & Harvested frozen from open models      & Specified \\
v3 & Mixture of \emph{Diverse} Experts    & Harvested frozen; hypermodal; native   & This paper; zen5-max \\
\bottomrule
\end{tabular}
\end{table}

\textbf{v1} (2026) applied Drop-Upcycling~\cite{dropupcycling}: clone a trained dense FFN into $N$
experts and inject structured noise to break the symmetry that otherwise collapses them back into
one, at roughly a quarter the cost of training an MoE from scratch. The insight it contributed ---
that expert \emph{diversity} must be engineered, because identical experts receive identical
gradients and never differentiate --- is the premise v2 inherited and took to its conclusion: if
diversity is what matters, the strongest available diversity is not noise injected into copies of
one model, but experts drawn from models that were independently trained by different groups, on
different data, with different attention mechanisms. Upcycling manufactures diversity; harvesting
finds it already made.

\textbf{v2} established that architecture --- harvest frozen, align, route by complexity ---
and specified it for zen5. It was never built: the reference implementation does not run, and
the zen5 SKU serves from an upstream model. No MoDE model exists yet.
\textbf{v3}, this paper, keeps v2's algebra and adds two things: placement as a first-class value
(Section~5) with the Map/Reduce/Route partition that follows from it (Section~6), and a hypermodal
router that selects encoders and decoders per modality rather than only FFN experts.

We record v1 rather than quietly drop it. An architecture line that erases its own supersessions
cannot be audited, and the reason v2 exists is legible only against what v1 taught. v1's technique
is not wrong either --- it is right wherever diversity has \emph{not} already been trained by
someone else, which is simply not the case for open frontier text models. Where that condition does
hold, upcycling is the correct method, and it is applied there~\cite{muen}.

MoDE ends here. What follows is not a fourth MoDE but a different line: the diffusion
work (enso) arrived independently at the idea of routing across experts too diverse to have been
trained together, and its own v3 generalizes past what MoDE binds. That architecture is
MUEN~\cite{muen}, and Section~6 is what MoDE contributes to it.


\section{Introduction}

Scaling a frontier model conventionally means training one: a single architecture, one
pretraining run, and a compute bill that grows with parameter count. Mixture-of-Experts
reduces the \emph{inference} cost of that model by activating a subset of its parameters,
but the experts are still trained together, from scratch, within one family.

The open-weight ecosystem now offers a different opportunity. Several organizations have
each spent enormous compute training large models with genuinely different inductive
biases---Gated DeltaNet, Lightning Attention, GLM MoE, DeepseekV3-style MoE, FP8 MoE.
Those weights are published. The expensive artifact---a trained expert---already exists;
what is missing is a way to use experts from different families \emph{together}.

MoDE is that mechanism. It treats published models as an \textbf{expert pool} to be
harvested, not a baseline to be beaten, and trains only the machinery needed to route
across them:

\begin{enumerate}
  \item \textbf{Diverse, not distilled.} Experts are taken \emph{as-is} from six model
        families and never fine-tuned. Their diversity is the point: models trained by
        different groups, on different data, with different attention mechanisms learn
        different representations, and a router can exploit that.
  \item \textbf{Complexity-aware routing.} A lightweight estimator predicts task
        difficulty from the first 64 tokens and selects a tier. Standard MoE routes every
        token through the same-size network; MoDE varies the network.
  \item \textbf{Frozen experts, tiny trainable surface.} $\sim$394M trained parameters
        against a $\sim$3.1T pool---two to three orders of magnitude less training compute
        than pretraining a model of comparable capability.
  \item \textbf{Pinned placement.} Frozen experts are pinned to devices. Residency is
        paid once; routing is then an address, not a schedule.
\end{enumerate}

\section{Expert Pool}

Experts are harvested from open-weight models spanning six text families and four
additional modalities. Each source contributes its FFN/MoE blocks; weights are frozen at
extraction and never updated.

\begin{table}[H]
\centering
\caption{Text expert sources. Tier is the complexity band the source serves.}
\begin{tabular}{llrrl}
\toprule
\textbf{Source} & \textbf{Architecture} & \textbf{Total} & \textbf{Active} & \textbf{Tier} \\
\midrule
Qwen3.5-0.8B & Gated DeltaNet (dense)       & 0.8B  & 0.8B & T0 --- trivial \\
Qwen3.5-9B   & Gated DeltaNet (dense)       & 9B    & 9B   & T1 --- standard \\
MiniMax-M2.5 & Lightning Attention MoE      & 230B  & 10B  & T2 --- complex \\
GLM-5        & GLM MoE                      & 744B  & 40B  & T3 --- advanced \\
Kimi K2.5    & DeepseekV3 MoE (384 experts) & 1.04T & 32B  & T4 --- frontier \\
Ling-1T      & FP8 MoE                      & 1T    & 50B  & T4 --- frontier \\
\bottomrule
\end{tabular}
\end{table}

\begin{table}[H]
\centering
\caption{Multimodal expert sources.}
\begin{tabular}{lllr}
\toprule
\textbf{Modality} & \textbf{Source} & \textbf{Architecture} & \textbf{Total} \\
\midrule
Vision & Qwen3-VL-30B          & Vision-language transformer   & 30B \\
Video  & Wan2.2                & MoE diffusion (3D causal VAE) & $\sim$14B \\
Video  & CogVideoX-5B          & Expert transformer + 3D-RoPE  & 5B \\
3D     & TRELLIS.2 (Microsoft) & Rectified flow DiT + SC-VAE   & 4B \\
Audio  & Qwen3-Omni            & Thinker--Talker omni MoE      & 30B \\
\bottomrule
\end{tabular}
\end{table}

The assembled pool is $\sim$3.1T total parameters across 1000+ experts, of which at most
$\sim$189B are active on any request. Every source is used under its own license and
credited, with that license, in \texttt{NOTICE}.

\section{Complexity-Aware Hierarchical Routing}

\subsection{The estimator}

Standard MoE applies the same computation to every token: a linear router selects top-$K$
experts from one pool, and depth is fixed whether the prompt is a greeting or a research
question. MoDE first asks how hard the task is.

A 207M-parameter estimator reads the first 64 tokens and predicts a tier
$t \in \{T0,\dots,T4\}$. Only that tier's experts are activated.

\begin{table}[H]
\centering
\caption{Complexity tiers and their activation budget.}
\begin{tabular}{llll}
\toprule
\textbf{Tier} & \textbf{Active} & \textbf{Latency target} & \textbf{Example} \\
\midrule
T0 & 0.8B      & $<$50ms  & ``Hello, how are you?'' \\
T1 & 9B        & $<$200ms & ``Summarize this article'' \\
T2 & 10--40B   & $<$2s    & ``Compare Keynesian vs.\ monetarist economics'' \\
T3 & 40--80B   & $<$5s    & ``Derive Black--Scholes from first principles'' \\
T4 & 80--100B+ & $<$10s   & ``Design a novel consensus algorithm'' \\
\bottomrule
\end{tabular}
\end{table}

The estimator is deliberately small: it must be cheap enough that running it on every
request is free relative to the compute it saves.

\subsection{Cross-architecture alignment}

Experts from different families do not share a representation space. A 134M-parameter
alignment layer projects each source's hidden states into a shared space, which is what
permits a Gated DeltaNet expert and a DeepseekV3 MoE expert to serve the same request.

\subsection{Router and escalation}

The 52M-parameter router performs per-tier expert selection, adopting two published
results rather than reinventing them: \textbf{ReMoE} ReLU routing (ICLR 2025) for an
adaptive expert count per token instead of a fixed top-$K$, and \textbf{MoE++}
zero-computation experts (ICLR 2025, oral)---zero, copy, and constant experts that let
trivial tokens bypass FFN compute, reported upstream at 1.1--2.1$\times$ throughput.

A tier chosen from 64 tokens can be wrong. When generation confidence drops below
threshold, MoDE promotes the request mid-generation to a higher tier. The estimator need
not be right, only cheap and approximately right; escalation makes its error recoverable
rather than fatal.

\section{Trainable Surface}

\begin{table}[H]
\centering
\caption{Trained parameters. Everything else is frozen.}
\begin{tabular}{lrl}
\toprule
\textbf{Component} & \textbf{Parameters} & \textbf{Role} \\
\midrule
ComplexityEstimator & 207M & Predicts the task's tier \\
AlignmentLayer      & 134M & Projects diverse architectures to a shared space \\
MoDERouter          & 52M  & Per-tier expert routing (ReLU gating) \\
\midrule
\textbf{Total} & \textbf{$\sim$394M} & \textbf{0.013\% of the $\sim$3.1T pool} \\
\bottomrule
\end{tabular}
\end{table}

This is the central economic claim, and it follows from the design rather than from a
measurement: if the experts are frozen and already trained, the training cost of the
assembled model is the training cost of the router.

\section{Placement: Pinned Experts}

A 3.1T pool does not fit on a device. The question is not whether to distribute it but
what distribution \emph{costs}, and here the frozen-ness of harvested experts pays a
second time.

\subsection{Frozen implies pinnable}

In conventional distributed training an expert is a liability: its weights change every
step, so it carries gradients and optimizer state, must be synchronized, and its
placement interacts with the training schedule. A MoDE expert changes never. That single
fact collapses the problem:

\begin{itemize}
  \item \textbf{No sync.} A frozen expert has no gradient and no optimizer state. Nothing
        about it needs to travel between devices, ever.
  \item \textbf{Residency is paid once.} Weights are loaded at extraction time and stay.
        There is no reload, no checkpoint restore into the hot path, no eviction pressure
        from an optimizer.
  \item \textbf{An expert is a pure function.} A frozen expert pinned to device $d$ is a
        deterministic map from an aligned hidden state to an aligned hidden state. It has
        no identity beyond its weights and no state beyond its input.
\end{itemize}

Placement is therefore a \textbf{value}, not a schedule: a static map
$\pi : \text{expert} \rightarrow \text{device}$, decided once and shared read-only. A
router that emits an expert id has, by composition with $\pi$, emitted an address. There
is no placement \emph{policy} at runtime to get wrong, because there is no decision left
to make.

\subsection{What pins where}

Tiering and placement compose. The T0/T1 experts are small dense models (0.8B, 9B); they
are cheap enough to replicate on \emph{every} device, so the overwhelmingly common case
---a trivial or standard request---is served with zero device hops. The large MoE sources
(T2--T4) are sharded across the fleet, and only the rare hard request pays a crossing.

\begin{table}[H]
\centering
\caption{Placement policy by tier. Replication is an optimization the frozen weights make free.}
\begin{tabular}{llll}
\toprule
\textbf{Tier} & \textbf{Source} & \textbf{Placement} & \textbf{Hops} \\
\midrule
T0/T1   & Qwen3.5-0.8B / 9B (dense) & Replicated per device & 0 \\
T2--T4  & MiniMax / GLM-5 / Kimi / Ling (MoE) & Sharded, pinned & 1 per crossing \\
Modal   & VL / video / 3D / audio & Pinned to modality hosts & 0--1 \\
\bottomrule
\end{tabular}
\end{table}

Because the distribution of real traffic is dominated by the cheap tiers, the expected
number of crossings per request is far below one---the fleet behaves like a small model
that occasionally consults a large one.

\subsection{Pinned and streamed are the same map}

Not every deployment has fleet-scale memory. A device may hold its pinned set resident
and stream the remainder from disk through an expert cache, which is viable precisely
because frozen weights are immutable and therefore trivially cacheable---a page of expert
weights is valid forever. Pinned and streamed are then not two systems but two
\emph{residency classes} under one placement map: $\pi$ says which device owns an expert;
residency says how fast that device can produce it.

\section{The Fusion Algebra: Map, Reduce, Route}

MoDE's serving requirements extend the kernel DSL (\texttt{hanzo-kernel}) by exactly one
class, and the extension is forced rather than designed.

\subsection{The existing law}

The DSL classes each operation by how it composes, and derives fusion from the class
alone---not from a pattern matcher over concrete operations:

\begin{itemize}
  \item \textbf{Map} --- index-local. $\text{map}(g) \circ \text{map}(f) = \text{map}(g
        \circ f)$. Map regions fuse into one kernel.
  \item \textbf{Reduce} --- row-local. Output depends on a whole row, so it is a
        \emph{fence}: fusion stops.
\end{itemize}

Fusion is decidable because the class, not the operation, decides it.

\subsection{The forced third class}

A MoDE expert is a Map region: an FFN chain (SwiGLU is $\text{silu}(xW_g) \odot xW_u$
followed by $W_d$) is pointwise between contractions, so within one expert the existing
algebra already fuses everything fusable. But an expert is Map \emph{on its own device}.
Routing a row to expert $e$ pinned at $\pi(e)$ is index-local in the same sense a Map is
---it depends on that row alone---yet it cannot fuse with its neighbours, because they may
compute somewhere else.

That is a fence, but not a contraction fence. It is a \textbf{placement fence}:

\begin{itemize}
  \item \textbf{Map} --- index-local, same device. Composes.
  \item \textbf{Reduce} --- row-local, same device. Fences \emph{within} a device.
  \item \textbf{Route} --- expert-local, possibly other device. Fences \emph{across}
        devices.
\end{itemize}

\lstset{caption={The class extension. Route is a fence like Reduce, distinguished by what it crosses.}}
\begin{lstlisting}
pub enum Class {
    Map,     // index-local  -> composes
    Reduce,  // row-local    -> fences (contraction)
    Route,   // expert-local -> fences (placement)
}
\end{lstlisting}

The consequence is that placement becomes decidable by the rule that already makes fusion
decidable. A trace is partitioned at its fences; \textsc{Reduce} fences split a kernel,
\textsc{Route} fences split a \emph{device}. Each resulting Map region is a fused kernel
that runs entirely on one device, and the device is $\pi(e)$ of the region's expert. The
scheduler does not need to understand experts, and the fuser does not need to understand
devices; both only read the class.

\subsection{Why this is the right seam}

The alternative---teaching the fuser about experts, or the router about kernels---braids
two concerns that have no reason to know each other. Classing \textsc{Route} keeps them
orthogonal: routing changes which region runs where, fusion changes what a region costs,
and neither changes the other's rules. Adding a seventh expert family adds rows to $\pi$
and nothing to the fuser. Adding a fusable op adds a variant to \texttt{UnOp} and nothing
to the router.

\subsection{Training under the same algebra}

Training uses the identical partition, which is what makes the trainable surface cheap to
train as well as small. The experts are frozen, so no \textsc{Route}-fenced region ever
requires a backward pass: gradients flow only through the alignment layer and the router,
both of which are Map/Reduce regions on the device that hosts them. The backward graph is
therefore a strict subgraph of the forward graph with every \textsc{Route} region cut out
---which is the algebraic statement of ``we train 0.013\% of the model.''

\section{Training}

\begin{enumerate}
  \item \textbf{Expert extraction} --- harvest FFN/MoE blocks from each source model.
  \item \textbf{Alignment pre-training} --- learn the cross-architecture projections.
  \item \textbf{Router training} --- train the complexity estimator and per-tier routers.
  \item \textbf{Integration} --- joint fine-tuning of the trainable surface, experts frozen.
  \item \textbf{Omnimodal fusion} --- add the vision, video, 3D, and audio experts.
\end{enumerate}

Hanzo's contribution on top of the harvested experts is confined to identity training,
agentic-data fine-tuning, and abliteration. The experts' capabilities are the source
models' own, and we claim none of them.

\section{Lineup}

\begin{table}[H]
\centering
\caption{Design targets, not shipped models. Active parameters are a range because the tier
is chosen per request. These SKU names currently serve from upstream open models; zen5-max is
intended to be the first that runs MoDE.}
\begin{tabular}{lrll}
\toprule
\textbf{Model} & \textbf{Total} & \textbf{Active} & \textbf{Target} \\
\midrule
zen5       & 750B & 0.8--50B   & General \\
zen5-coder & 1.8T & 0.8--80B   & Code \\
zen5-omni  & 2.5T & 0.8--100B  & Omnimodal \\
zen5-max   & 3.1T & 0.8--100B+ & Frontier \\
\bottomrule
\end{tabular}
\end{table}

\section{Evaluation}

There is no evaluation, because there is nothing yet to evaluate. Routing accuracy,
per-tier latency, crossing rate under real traffic, and end-task quality require a
harness that does not exist: the reference implementation in the zen5 repository has
never been executed, and we do not cite it as one. We publish no numbers here.

Two claims decide whether the architecture holds---frontier capability at
router-training cost, and sub-unity expected crossings per request---and both rest on a
premise that has itself never been tested: that a trained projection makes another
family's frozen FFN useful (Section~3.2). If that premise fails, MoDE as specified does
not work, and no amount of the rest repairs it. It is the first thing to measure and the
cheapest to falsify. A fabricated table would only obscure the one question worth asking.

\section{Related Work}

MoDE builds directly on published results: Drop-Upcycling~\cite{dropupcycling}, which v1 was
built on and whose diagnosis of weight-correlation collapse motivates harvesting; ReMoE
(ICLR 2025) for ReLU-gated adaptive expert counts; MoE++ (ICLR 2025, oral) for zero-computation
experts; and Transfusion for unified autoregressive and diffusion modeling in one
architecture. Model merging and
FrankenMoE assemblies also combine published checkpoints, but they merge weights into a
single homogeneous model; MoDE keeps the sources distinct, aligned rather than averaged,
and selects among them per request. Its novelty is the composition: cross-architecture
experts harvested from independently trained open models, made interoperable by a learned
alignment, selected by a complexity estimator that varies compute with task difficulty,
and pinned so that selection and placement are the same decision. We are not aware of an
existing model that combines these.

\section*{Availability}

Architecture, extraction, and router-training code: \url{https://github.com/zenlm/zen5}
(a reference implementation; it does not run, and the native one is being built on the Rust
stack).
Kernel DSL and the Map/Reduce/Route partition: \url{https://github.com/hanzoai/ml}
(\texttt{hanzo-kernel}). The successor line, MUEN: \url{https://github.com/hanzoai/muen}. Each
harvested source is credited, with its license, in \texttt{NOTICE}.

\begin{thebibliography}{9}

\bibitem{muen}
Zen LM Research, Hanzo AI.
\emph{MUEN: Multimodal Mixture of Unbound Experts}. enso v3 (DiT-MoDE).
\url{https://github.com/hanzoai/muen}

\bibitem{dropupcycling}
Drop-Upcycling: Training Sparse Mixture of Experts with Partial Re-initialization.
\url{https://arxiv.org/abs/2502.19261}
\end{thebibliography}

\end{document}
